\clearpage
\chapter{Conclusion}
\section{Conclusion}
This project addressed the inconsistency of slice-wise automated glioma segmentation by adapting the cross-slice transformer idea from CT lymph-node segmentation to multi-modal brain MRI. We designed and implemented a 2.5D Cross-Slice U-Net that encodes five neighbouring slices with a shared encoder and fuses them through self-attention across the slice axis at the bottleneck. The same code reduces to a 2D U-Net baseline when $K=1$. We replaced the mid-review prototype's evaluation and data handling with a rigorous pipeline: original NIfTI data normalised per patient, patient-level five-fold cross-validation with separate validation and test sets, training on all slices, and volumetric per-patient metrics with paired statistical tests. We also justified the hyper-parameters with a design study and an Optuna search, and ran the full-scale experiments on the institutional Astra A100 cluster.

The objectives were achieved as follows. The model was implemented (objective 1), and the leakage-free, reproducible pipeline is publicly available (objective 2). The controlled five-fold comparison on all 369 BraTS 2020 patients showed that cross-slice context improves Dice for all three regions (WT 0.904 vs 0.898, TC 0.840 vs 0.837, ET 0.789 vs 0.783), with small but statistically significant gains (objective 3). The hyper-parameters were justified, and the Optuna search independently confirmed the chosen configuration (objective 4). Boundary error and slice consistency improved markedly: HD95 fell by 1.8~mm for WT and 5.1~mm for ET, and false-positive slices fell by 34\% (objective 5). The full experiments were completed on the HPC cluster (objective 6). These findings are consistent with the results discussed in Chapter~6 and support the conclusion that cross-slice attention is a compact and effective way to make 2.5D glioma segmentation more consistent through the volume. In the context of SDG 3, the demonstrated contribution is a lightweight model with better volumetric consistency; the reduction of clinicians' workload and diagnostic delay remains an expected benefit that requires clinical evaluation.

\section{Future work}
Several tasks remain for the next phase. The most immediate is to quantify robustness by retraining both models with several random seeds per fold, and to submit predictions to the official BraTS validation server so that results become directly comparable with published methods. A qualitative evaluation, with visual overlays of predictions and review by a radiologist, is needed to confirm that the measured gains in boundary accuracy are clinically meaningful. A 3D U-Net or nnU-Net reference trained under the same splits would show how much of the gap between 2.5D and 3D models the cross-slice module closes, and at what computational cost. The proposed data-size study, training both models on 16, 64, 150 and 265 patients with the corrected protocol, would test whether cross-slice context helps more when labelled data are scarce, as the preliminary mid-review study suggested; this is relevant for hospitals with few annotated scans. On the modelling side, the descriptor-based fusion of the original paper should be implemented as a second variant, and attention could be applied at the skip-connection levels as well as the bottleneck. Finally, evaluation on external and more recent datasets (such as later BraTS editions), robustness to missing modalities, and calibrated uncertainty estimates would be prerequisites for any clinical use.
